\unitchapter{2}{4}{Database Management Systems}{p2-04}

\unitmeta{Expected questions: 10--12 · Target: 10+ · Type: normalisation, keys,
SQL output, transactions, concepts --- very scoring}

\begin{sylbox}

\begin{itemize}
\tightlist
\item[$\square$]
  Database system concepts \& architecture: data models, schemas, instances, three-schema architecture, data independence, database languages \& interfaces, centralised \& client/server architecture
\item[$\square$]
  Data modelling: ER diagram, relational model (constraints, languages, design, programming), relational algebra \& calculus, codd rules
\item[$\square$]
  SQL: data definition, types, constraints, queries, insert, delete, update, views, stored procedures, functions, triggers, SQL injection
\item[$\square$]
  Normalisation: functional dependencies, 1NF, 2NF, 3NF, BCNF, multivalued dependencies \& 4NF, join dependencies \& 5NF, inclusion dependencies
\item[$\square$]
  Transaction processing, concurrency control \& recovery
\item[$\square$]
  Physical database design: file organisation, indexing (single-level, multilevel, B-tree, B+-tree), query processing \& optimisation
\item[$\square$]
  Object \& object-relational databases; database security
\item[$\square$]
  Enhanced data models: temporal, multimedia, deductive, XML \& internet databases, mobile, GIS, genome, distributed databases
\item[$\square$]
  Data warehousing \& data mining: OLAP, data cube, association rules, classification, clustering
\item[$\square$]
  Big data systems: characteristics, types, architecture, MapReduce, Hadoop, NoSQL (CAP theorem, document, key-value, column, graph)
\end{itemize}

\end{sylbox}

\needspace{161pt}

\hypertarget{p2-04--1-architecture}{%
\section{Architecture}\label{p2-04--1-architecture}}

A database management system stands between users and stored data. Its central achievement is \emph{data
independence}: applications describe what data they want, and the system decides how it is stored and
found, so that storage can change without rewriting programs.

\needspace{95pt}

\hypertarget{p2-04--11-three-schema-ansisparc-architecture}{%
\subsection{Three-Schema (ANSI/SPARC) Architecture}\label{p2-04--11-three-schema-ansisparc-architecture}}

\begin{figure}[!htbp]
\centering
\begin{adjustbox}{max width=\linewidth}
\begin{tikzpicture}
  \fill[fslate!35] (-1,2.5) rectangle (11,4.1); \fill[fsage!45] (-1,0.9) rectangle (11,2.3); \fill[fsand!55] (-1,-0.7) rectangle (11,0.7);
  \node[capt, anchor=west] at (-0.9,3.85) {external level}; \node[capt, anchor=west] at (-0.9,2.05) {conceptual level}; \node[capt, anchor=west] at (-0.9,0.45) {internal level};
  \foreach \x/\n in {3/1, 5.4/2, 7.8/3} {
    \node[box, fill=white, minimum width=20mm, font=\footnotesize] (v\n) at (\x,3.2) {external view \n};
    \node[font=\scriptsize] (u\n) at (\x,4.55) {user \n};
    \draw[arr] (u\n) -- (v\n);
  }
  \node[box, fill=white, text width=58mm, font=\footnotesize] (C) at (5.4,1.6) {\textbf{conceptual schema} — entities, relationships, constraints};
  \node[box, fill=white, text width=58mm, font=\footnotesize] (I) at (5.4,0.0) {\textbf{internal schema} — storage structures, indexes, files};
  \node[db, minimum width=26mm, minimum height=13mm] (D) at (5.4,-1.8) {stored\\database};
  \foreach \n in {1,2,3} \draw[arrd] (v\n) -- (C);
  \draw[arrd] (C) -- (I); \draw[arrd] (I) -- (D);
  \node[lbl, anchor=west] at (8.8,2.4) {logical data independence};
  \node[lbl, anchor=west] at (8.4,0.8) {physical data independence};
\end{tikzpicture}
\end{adjustbox}
\caption{The ANSI/SPARC three-schema architecture. Mappings between adjacent levels insulate each level from
changes in the one below it.}
\end{figure}


\begin{itemize}
\tightlist
\item
  \textbf{Logical data independence}: change conceptual schema without changing external views (harder to achieve).
\item
  \textbf{Physical data independence}: change internal schema without changing conceptual schema.
\item
  \textbf{Schema} = structure (intension, changes rarely); \textbf{Instance} = data at a moment (extension).
\item
  Languages: \textbf{DDL} (CREATE, ALTER, DROP, TRUNCATE), \textbf{DML} (SELECT, INSERT, UPDATE, DELETE), \textbf{DCL} (GRANT, REVOKE), \textbf{TCL} (COMMIT, ROLLBACK, SAVEPOINT).
\item
  Data models: hierarchical (IMS, tree), network (CODASYL, graph), relational (Codd 1970), object-oriented, object-relational.
\item
  Architectures: centralised, \textbf{2-tier} client/server, \textbf{3-tier} (client -- application server -- DB server).
\end{itemize}

\needspace{191pt}

\hypertarget{p2-04--12-dbms-components}{%
\subsection{DBMS Components}\label{p2-04--12-dbms-components}}

Query processor (DDL interpreter, DML compiler, query evaluation engine), storage manager (authorisation, transaction manager, file manager, buffer manager), data dictionary (metadata).

\needspace{125pt}

\hypertarget{p2-04--2-er-model}{%
\section{ER Model}\label{p2-04--2-er-model}}

The entity--relationship model is a picture of the real world before any tables exist: things (entities),
their properties (attributes) and the associations between them (relationships). Chen\textquotesingle s notation, shown
below, is the one used in the examination.

\begin{figure}[!htbp]
\centering
\begin{adjustbox}{max width=\linewidth}
\begin{tikzpicture}[ent/.style={ink, rectangle, fill=fslate, minimum width=19mm, minimum height=7mm, font=\footnotesize},
    att/.style={ink, ellipse, fill=fsand, minimum width=17mm, minimum height=7mm, font=\footnotesize, inner sep=1pt},
    rel/.style={ink, diamond, fill=fsage, aspect=1.8, inner sep=1pt, minimum width=17mm, minimum height=9mm, font=\footnotesize}]
  \node[ent] at (0,0) {Entity}; \node[lbl] at (0,-0.75) {entity set};
  \node[ent, double, double distance=1.5pt] at (3.4,0) {Weak}; \node[lbl] at (3.4,-0.75) {weak entity};
  \node[att] at (6.8,0) {attribute}; \node[lbl] at (6.8,-0.75) {simple attribute};
  \node[att] at (10.2,0) {\underline{roll\_no}}; \node[lbl] at (10.2,-0.75) {key attribute};
  \node[att, double, double distance=1.3pt] at (0,-2.0) {phones}; \node[lbl] at (0,-2.75) {multivalued};
  \node[att, dashed] at (3.4,-2.0) {age}; \node[lbl] at (3.4,-2.75) {derived};
  \node[rel] at (6.8,-2.0) {Rel}; \node[lbl] at (6.8,-2.85) {relationship};
  \node[rel, double, double distance=1.3pt] at (10.2,-2.0) {Ident}; \node[lbl] at (10.2,-2.85) {identifying relationship};
  \draw[fgink, line width=0.5pt] (1.0,-4.0) -- (3.6,-4.0); \node[lbl, anchor=west] at (3.8,-4.0) {partial participation};
  \draw[fgink, line width=0.5pt, double, double distance=1.4pt] (7.2,-4.0) -- (9.8,-4.0); \node[lbl, anchor=west] at (10.0,-4.0) {total participation};
\end{tikzpicture}
\end{adjustbox}
\caption{Chen's ER notation.}
\end{figure}


\begin{figure}[!htbp]
\centering
\begin{adjustbox}{max width=\linewidth}
\begin{tikzpicture}[ent/.style={ink, rectangle, fill=fslate, minimum width=22mm, minimum height=8mm, font=\small\bfseries},
    att/.style={ink, ellipse, fill=fsand, minimum height=6.5mm, font=\scriptsize, inner sep=1.5pt},
    rel/.style={ink, diamond, fill=fsage, aspect=1.9, inner sep=1pt, font=\footnotesize\bfseries}]
  \node[ent] (S) at (0,0) {STUDENT};
  \node[rel] (E) at (3.6,0) {ENROLLS};
  \node[ent] (C) at (7.2,0) {COURSE};
  \node[rel] (O) at (10.6,0) {OFFERS};
  \node[ent] (D) at (13.6,0) {DEPARTMENT};
  \draw[thinline] (S) -- node[above, font=\scriptsize] {M} (E);
  \draw[thinline] (E) -- node[above, font=\scriptsize] {N} (C);
  \draw[thinline, double, double distance=1.3pt] (C) -- node[above, font=\scriptsize] {N} (O);
  \draw[thinline] (O) -- node[above, font=\scriptsize] {1} (D);
  \node[att] (s1) at (-1.3,1.4) {\underline{roll\_no}}; \node[att] (s2) at (0.2,1.6) {name}; \node[att] (s3) at (1.5,1.3) {dob};
  \foreach \a in {s1,s2,s3} \draw[thinline] (S) -- (\a);
  \node[att] (e1) at (3.6,1.4) {grade}; \draw[thinline] (E) -- (e1);
  \node[att] (c1) at (6.0,1.4) {\underline{course\_id}}; \node[att] (c2) at (7.4,1.6) {title}; \node[att] (c3) at (8.6,1.3) {credits};
  \foreach \a in {c1,c2,c3} \draw[thinline] (C) -- (\a);
  \node[att] (d1) at (13.6,1.4) {\underline{dept\_id}}; \draw[thinline] (D) -- (d1);
\end{tikzpicture}
\end{adjustbox}
\caption{An ER diagram for a university. ENROLLS is many-to-many and carries the attribute \emph{grade};
every course must be offered by a department (total participation, double line).}
\end{figure}


\begin{itemize}
\tightlist
\item
  \textbf{Cardinality}: 1:1, 1:N, M:N. \textbf{Participation}: total (every entity participates --- double line) or partial.
\item
  \textbf{Weak entity}: no key of its own; identified via owner entity + partial key (discriminator).
\item
  \textbf{Specialisation} (top-down), \textbf{generalisation} (bottom-up), \textbf{aggregation} (relationship as entity). Disjoint/overlapping, total/partial constraints.
\end{itemize}

\needspace{130pt}

\hypertarget{p2-04--er--relational-mapping-minimum-tables}{%
\subsection{ER → Relational mapping (minimum tables)}\label{p2-04--er--relational-mapping-minimum-tables}}

\begingroup\small

\par\medskip\noindent\hfil\begin{tabular}{@{}
  >{\raggedright\arraybackslash}p{(\columnwidth - 2\tabcolsep) * \real{0.1917}}
  >{\raggedright\arraybackslash}p{(\columnwidth - 2\tabcolsep) * \real{0.4928}}@{}}
\toprule\noalign{}
\begin{minipage}[b]{\linewidth}\raggedright
\textbf{Case}
\end{minipage} & \begin{minipage}[b]{\linewidth}\raggedright
\textbf{Tables}
\end{minipage} \\
\midrule\noalign{}
Strong entity & 1 table each \\
1:1 relationship & Merge into either side (total participation side preferred) \\
1:N relationship & FK on N-side; no separate table \\
\textbf{M:N relationship} & \textbf{Separate table} with both PKs \\
Multivalued attribute & Separate table \\
Weak entity & Table with owner\textquotesingle s PK + partial key \\
\bottomrule\noalign{}
\end{tabular}\hfil\null\par\medskip


\endgroup

Classic: E1 ---M:N--- E2 ---1:N--- E3 → minimum tables = 3 entities + 1 (for M:N) = \textbf{4}.

\needspace{125pt}

\hypertarget{p2-04--3-keys--integrity}{%
\section{Keys \& Integrity}\label{p2-04--3-keys--integrity}}

Keys are how a relational database identifies rows and links tables. The three key concepts are nested, as
the figure shows; integrity constraints then forbid NULL primary keys and dangling foreign keys.

\begin{figure}[!htbp]
\centering
\begin{adjustbox}{max width=\linewidth}
\begin{tikzpicture}
  \filldraw[fill=fgrey, draw=fgink, line width=0.5pt] (0,0) ellipse (4.2 and 2.0);
  \filldraw[fill=fslate, draw=fgink, line width=0.5pt] (0.7,-0.2) ellipse (2.7 and 1.35);
  \filldraw[fill=fsand, draw=fgink, line width=0.5pt] (1.6,-0.35) ellipse (1.1 and 0.62);
  \node[font=\small, align=center] at (-2.75,0.55) {\textbf{super keys}\\\footnotesize any set that\\\footnotesize identifies a tuple};
  \node[font=\small, align=center] at (-0.4,0.15) {\textbf{candidate}\\\textbf{keys}\\\footnotesize minimal};
  \node[font=\small, align=center] at (1.6,-0.35) {\textbf{primary}\\\footnotesize not NULL};
  \node[font=\footnotesize, align=left, anchor=west] at (4.5,0.4) {alternate keys: candidate keys\\not chosen as primary};
  \node[font=\footnotesize, align=left, anchor=west] at (4.5,-0.7) {foreign key: refers to the\\primary key of another relation};
\end{tikzpicture}
\end{adjustbox}
\caption{Every primary key is a candidate key, and every candidate key is a (minimal) super key.}
\end{figure}


\begin{itemize}
\tightlist
\item
  Number of super keys of R(A1\ldots An) with single candidate key A1: \textbf{2ⁿ⁻¹}.
\item
  With candidate keys A1 and A2: 2ⁿ⁻¹ + 2ⁿ⁻¹ − 2ⁿ⁻² = \textbf{3·2ⁿ⁻²}.
\item
  \textbf{Entity integrity}: PK cannot be NULL. \textbf{Referential integrity}: FK value must match an existing PK or be NULL.
\item
  ON DELETE CASCADE / SET NULL / RESTRICT.
\end{itemize}

\textbf{Codd\textquotesingle s 12 rules (13 with Rule 0)}: Rule 0 foundation; 1 information; 2 guaranteed access; 3 systematic NULLs; 4 active online catalog; 5 comprehensive data sublanguage; 6 view updating; 7 high-level insert/update/delete; 8 physical data independence; 9 logical data independence; 10 integrity independence; 11 distribution independence; 12 non-subversion.

\needspace{160pt}

\hypertarget{p2-04--4-relational-algebra--calculus}{%
\section{Relational Algebra \& Calculus}\label{p2-04--4-relational-algebra--calculus}}

Relational algebra is a small set of operations on tables that produce tables. Every SQL query is
translated into an algebra expression before it is optimised, which is why the algebra is examined so
often.

\begingroup\small

\par\medskip\noindent\hfil\begin{tabular}{@{}
  >{\raggedright\arraybackslash}p{(\columnwidth - 4\tabcolsep) * \real{0.1684}}
  >{\raggedright\arraybackslash}p{(\columnwidth - 4\tabcolsep) * \real{0.2178}}
  >{\raggedright\arraybackslash}p{(\columnwidth - 4\tabcolsep) * \real{0.2787}}@{}}
\toprule\noalign{}
\begin{minipage}[b]{\linewidth}\raggedright
\textbf{Operation}
\end{minipage} & \begin{minipage}[b]{\linewidth}\raggedright
\textbf{Symbol}
\end{minipage} & \begin{minipage}[b]{\linewidth}\raggedright
\textbf{Notes}
\end{minipage} \\
\midrule\noalign{}
Selection & σ\_cond(R) & Rows; commutative \\
Projection & π\_attrs(R) & Columns; removes duplicates \\
Union & R ∪ S & Union-compatible \\
Difference & R − S & \\
Cartesian product & R × S & \\
Rename & ρ & \\
Intersection & R ∩ S = R − (R − S) & Derived \\
Natural join & R ⋈ S & Equal common attrs \\
Theta join & R ⋈\_θ S & σ\_θ(R × S) \\
Division & R ÷ S & "for all" queries \\
Outer joins & ⟕ ⟖ ⟗ & Keep unmatched with NULLs \\
\bottomrule\noalign{}
\end{tabular}\hfil\null\par\medskip


\endgroup

Fundamental: \textbf{σ, π, ∪, −, ×, ρ}. Relational algebra is \textbf{procedural}; relational calculus (TRC/DRC) is \textbf{non-procedural}. Safe calculus = relational algebra in power (relational completeness).

\begin{center}
\renewcommand{\arraystretch}{1.3}
\begin{tabular}{l c}
\toprule
\textbf{Operation} ($|R| = m$, $|S| = n$) & \textbf{Number of tuples} \\
\midrule
$R \times S$ & $m\,n$ \\
$R \bowtie S$ & $0 \dots m\,n$ \\
$R \cup S$ & $\max(m, n) \dots m + n$ \\
$R \cap S$ & $0 \dots \min(m, n)$ \\
$R - S$ & $0 \dots m$ \\
$R$ left outer join $S$ & at least $m$ \\
\bottomrule
\end{tabular}
\end{center}


Division example: "students who took \textbf{all} courses": π\_\{sid,cid\}(Enrolled) ÷ π\_\{cid\}(Course).

TRC: \texttt{\{\ t\ \textbar{}\ t\ ∈\ Student\ ∧\ t.age\ \textgreater{}\ 20\ \}} · DRC: \texttt{\{\ \textless{}n\textgreater{}\ \textbar{}\ ∃a\ (\textless{}n,a\textgreater{}\ ∈\ Student\ ∧\ a\ \textgreater{}\ 20)\ \}}.

\needspace{136pt}

\hypertarget{p2-04--5-sql}{%
\section{SQL}\label{p2-04--5-sql}}

\begin{listingblock}{86pt}
\begin{Verbatim}[fontsize=\fontsize{9.0}{10.9}\selectfont]
SELECT dept, COUNT(*) AS n, AVG(salary)
FROM   employee
WHERE  salary > 20000          -- filters rows (before grouping)
GROUP  BY dept
HAVING COUNT(*) > 2            -- filters groups
ORDER  BY n DESC;
\end{Verbatim}
\end{listingblock}

Logical order of execution: \textbf{FROM → WHERE → GROUP BY → HAVING → SELECT → DISTINCT → ORDER BY → LIMIT}.

\begin{itemize}
\tightlist
\item
  \texttt{COUNT(*)} counts all rows; \texttt{COUNT(col)} ignores NULLs; aggregates (SUM/AVG/MAX/MIN) ignore NULLs.
\item
  NULL comparisons give UNKNOWN; use \texttt{IS\ NULL}. \texttt{NULL\ =\ NULL} → UNKNOWN.
\item
  \texttt{DELETE} (DML, row-wise, rollback possible, WHERE allowed) vs \texttt{TRUNCATE} (DDL, all rows, faster) vs \texttt{DROP} (removes table structure).
\item
  \texttt{IN}, \texttt{EXISTS}, \texttt{ANY/SOME}, \texttt{ALL}: \texttt{x\ \textgreater{}\ ALL\ (subquery)} = greater than max; \texttt{x\ \textgreater{}\ ANY} = greater than min.
\item
  \texttt{UNION} removes duplicates; \texttt{UNION\ ALL} keeps them.
\item
  Correlated subquery: inner query refers to outer --- evaluated per outer row.
\item
  Constraints: PRIMARY KEY, FOREIGN KEY, UNIQUE (allows NULL), NOT NULL, CHECK, DEFAULT.
\item
  \textbf{View}: virtual table (stored query); updatable if based on single table without aggregates/DISTINCT/GROUP BY. \textbf{Materialised view} stores results.
\item
  \textbf{Trigger}: ECA (Event--Condition--Action) rule; BEFORE/AFTER, row-level/statement-level.
\item
  \textbf{Stored procedure} (CALL, may not return value) vs \textbf{function} (returns value, usable in SELECT).
\item
  \textbf{SQL injection}: attacker inserts SQL via input, e.g. \texttt{\textquotesingle{}\ OR\ \textquotesingle{}1\textquotesingle{}=\textquotesingle{}1}. Prevent with \textbf{prepared statements/parameterised queries}, input validation, least privilege.
\end{itemize}

\textbf{Nth highest salary}:

\begin{listingblock}{31pt}
\begin{Verbatim}[fontsize=\fontsize{9.0}{10.9}\selectfont]
SELECT MAX(salary) FROM emp WHERE salary < (SELECT MAX(salary) FROM emp);   -- 2nd highest
\end{Verbatim}
\end{listingblock}

\needspace{161pt}

\hypertarget{p2-04--6-functional-dependencies--normalisation}{%
\section{Functional Dependencies \& Normalisation}\label{p2-04--6-functional-dependencies--normalisation}}

A functional dependency \(X \to Y\) says that the value of \(X\) determines the value of \(Y\). Redundancy arises
when a table stores such facts about only part of its key, or about non-key attributes; normalisation
removes it by splitting the table so that every fact is stored once.

\needspace{95pt}

\hypertarget{p2-04--61-armstrongs-axioms}{%
\subsection{Armstrong\textquotesingle s Axioms}\label{p2-04--61-armstrongs-axioms}}

\begin{itemize}
\tightlist
\item
  \textbf{Reflexivity}: Y ⊆ X ⇒ X → Y · \textbf{Augmentation}: X → Y ⇒ XZ → YZ · \textbf{Transitivity}: X → Y, Y → Z ⇒ X → Z
\item
  Derived: Union, Decomposition, Pseudo-transitivity (X → Y, WY → Z ⇒ WX → Z).
\item
  Sound and complete.
\end{itemize}

\needspace{125pt}

\hypertarget{p2-04--62-attribute-closure--finding-candidate-keys}{%
\subsection{Attribute Closure --- Finding Candidate Keys}\label{p2-04--62-attribute-closure--finding-candidate-keys}}

R(A, B, C, D, E), F = \{A → B, B → C, CD → E\}

\begin{formulas}
\begin{align*}
A^{+} &= \{A, B, C\}\\
(AD)^{+} &= \{A, D, B, C, E\} = \text{all attributes} \;\Rightarrow\; AD \text{ is a key}\\
\end{align*}
\end{formulas}
\noindent Since $A$ and $D$ never appear on a right-hand side, every key must contain both; $AD$ is therefore the only candidate key.



Trick: attributes not on any RHS must be in every key; attributes only on RHS are never in a key.

\textbf{Minimal (canonical) cover}: single attribute on RHS → remove extraneous LHS attributes → remove redundant FDs.

\needspace{125pt}

\hypertarget{p2-04--63-normal-forms}{%
\subsection{Normal Forms}\label{p2-04--63-normal-forms}}

Each normal form forbids one more kind of undesirable dependency. The nested diagram is a reminder that
higher normal forms imply all the lower ones.

\begin{figure}[!htbp]
\centering
\begin{adjustbox}{max width=\linewidth}
\begin{tikzpicture}
  \foreach \w/\col/\name/\cond [count=\k] in {14.0/fgrey/1NF/{atomic values}, 12.0/fslate/2NF/{no partial dependency}, 10.0/fsage/3NF/{no transitive dependency},
      8.0/fsand/BCNF/{every determinant a super key}, 6.0/frose/4NF/{no non-trivial MVD}, 4.0/fochre/5NF/{no non-trivial join dependency}} {
    \pgfmathsetmacro{\yb}{0.25*(\k-1)}
    \pgfmathsetmacro{\yt}{5.6-0.6*(\k-1)}
    \filldraw[fill=\col, draw=fgink, line width=0.5pt, rounded corners=4pt] ({-\w/2},\yb) rectangle ({\w/2},\yt);
    \ifnum\k<6
      \node[font=\small, anchor=north] at (0,{\yt-0.06}) {\textbf{\name}\ \ \textit{\footnotesize \cond}};
    \else
      \node[font=\small, align=center] at (0,{(\yb+\yt)/2}) {\textbf{\name}\\\textit{\footnotesize \cond}};
    \fi
  }
\end{tikzpicture}
\end{adjustbox}
\caption{Each normal form is stricter than the one enclosing it: a relation in BCNF is also in 3NF, 2NF and 1NF.}
\end{figure}


\begingroup\small

\par\medskip\noindent\hfil\begin{tabular}{@{}
  >{\raggedright\arraybackslash}p{(\columnwidth - 2\tabcolsep) * \real{0.0767}}
  >{\raggedright\arraybackslash}p{(\columnwidth - 2\tabcolsep) * \real{0.7217}}@{}}
\toprule\noalign{}
\begin{minipage}[b]{\linewidth}\raggedright
\textbf{NF}
\end{minipage} & \begin{minipage}[b]{\linewidth}\raggedright
\textbf{Condition (for every non-trivial FD X → A)}
\end{minipage} \\
\midrule\noalign{}
1NF & All attributes atomic (no repeating groups / multivalued) \\
2NF & 1NF + no \textbf{partial} dependency (non-prime attribute depends on part of a candidate key) \\
3NF & X is a super key \textbf{OR} A is a \textbf{prime} attribute \\
BCNF & X is a \textbf{super key} \\
4NF & For every non-trivial MVD X →→ Y, X is a super key \\
5NF & Every join dependency is implied by candidate keys \\
\bottomrule\noalign{}
\end{tabular}\hfil\null\par\medskip


\endgroup

\begin{itemize}
\tightlist
\item
  Relation with \textbf{only two attributes is always in BCNF}.
\item
  If all candidate keys are single attributes ⇒ automatically 2NF.
\item
  If all attributes are prime ⇒ automatically 3NF.
\end{itemize}

\needspace{130pt}

\hypertarget{p2-04--64-decomposition-properties}{%
\subsection{Decomposition Properties}\label{p2-04--64-decomposition-properties}}

\begingroup\small

\par\medskip\noindent\hfil\begin{tabular}{@{}
  >{\raggedright\arraybackslash}p{(\columnwidth - 2\tabcolsep) * \real{0.3006}}
  >{\raggedright\arraybackslash}p{(\columnwidth - 2\tabcolsep) * \real{0.6006}}@{}}
\toprule\noalign{}
\begin{minipage}[b]{\linewidth}\raggedright
\textbf{Property}
\end{minipage} & \begin{minipage}[b]{\linewidth}\raggedright
\textbf{Test}
\end{minipage} \\
\midrule\noalign{}
\textbf{Lossless join} (must) & R into R1, R2 is lossless iff (R1 ∩ R2) → R1 or (R1 ∩ R2) → R2 \\
\textbf{Dependency preserving} (desirable) & (F1 ∪ F2)⁺ = F⁺ \\
\bottomrule\noalign{}
\end{tabular}\hfil\null\par\medskip


\endgroup

\begin{itemize}
\tightlist
\item
  \textbf{3NF} decomposition: always lossless \textbf{and} dependency-preserving (synthesis algorithm).
\item
  \textbf{BCNF} decomposition: always lossless but \textbf{not always} dependency-preserving.
\end{itemize}

\textbf{Worked}: R(A,B,C), F = \{AB → C, C → B\}. Keys: AB, AC. C → B: C not super key but B is prime → \textbf{3NF, not BCNF}. BCNF decomposition (C,B), (A,C) loses AB → C.

\needspace{196pt}

\hypertarget{p2-04--7-transactions}{%
\section{Transactions}\label{p2-04--7-transactions}}

A transaction is a unit of work --- such as a fund transfer --- that must happen completely or not at all, even
when many transactions run at once and the system may crash. The ACID properties state these guarantees.

\needspace{130pt}

\hypertarget{p2-04--acid}{%
\subsection{ACID}\label{p2-04--acid}}

\begingroup\small

\par\medskip\noindent\hfil\begin{tabular}{@{}
  >{\raggedright\arraybackslash}p{(\columnwidth - 2\tabcolsep) * \real{0.2372}}
  >{\raggedright\arraybackslash}p{(\columnwidth - 2\tabcolsep) * \real{0.3100}}@{}}
\toprule\noalign{}
\begin{minipage}[b]{\linewidth}\raggedright
\textbf{Property}
\end{minipage} & \begin{minipage}[b]{\linewidth}\raggedright
\textbf{Ensured by}
\end{minipage} \\
\midrule\noalign{}
\textbf{A}tomicity (all or nothing) & Recovery manager (undo log) \\
\textbf{C}onsistency & Programmer + integrity constraints \\
\textbf{I}solation & Concurrency control manager \\
\textbf{D}urability & Recovery manager (redo log) \\
\bottomrule\noalign{}
\end{tabular}\hfil\null\par\medskip


\endgroup

\needspace{95pt}

\hypertarget{p2-04--transaction-states}{%
\subsection{Transaction States}\label{p2-04--transaction-states}}

\begin{figure}[!htbp]
\centering
\begin{adjustbox}{max width=\linewidth}
\begin{tikzpicture}[st/.style={ink, rounded rectangle, fill=fslate, minimum width=24mm, minimum height=8mm, font=\small,
    blur shadow={shadow blur steps=6, shadow xshift=0.7pt, shadow yshift=-0.9pt, shadow opacity=22}}]
  \node[st] (A) at (0,0) {active};
  \node[st] (P) at (4,1.6) {partially committed};
  \node[st, fill=fsage] (C) at (9.6,1.6) {committed};
  \node[st, fill=frose] (F) at (4,-1.6) {failed};
  \node[st, fill=fgrey] (B) at (9.6,-1.6) {aborted};
  \draw[arr] (-1.6,0) -- (A);
  \draw[arr] (A) -- node[lbl, sloped, above] {last statement} (P);
  \draw[arr] (P) -- node[lbl, above] {written to stable storage} (C);
  \draw[arr] (A) -- node[lbl, sloped, below] {error} (F);
  \draw[arr] (P) -- node[lbl, right] {error} (F);
  \draw[arr] (F) -- node[lbl, below] {rolled back} (B);
\end{tikzpicture}
\end{adjustbox}
\caption{The states of a transaction.}
\end{figure}


\needspace{125pt}

\hypertarget{p2-04--schedules--serializability}{%
\subsection{Schedules \& Serializability}\label{p2-04--schedules--serializability}}

When transactions interleave, the result is correct if it is equivalent to some serial order. Conflict
serializability is tested by drawing a precedence graph and looking for a cycle.

\begin{itemize}
\tightlist
\item
  \textbf{Conflicting operations}: same data item, different transactions, at least one write (RW, WR, WW).
\item
  \textbf{Conflict serializable} ⇔ \textbf{precedence graph is acyclic} (edge Ti → Tj if Ti\textquotesingle s op conflicts with and precedes Tj\textquotesingle s).
\item
  \textbf{View serializable}: same initial reads, same reads-from, same final writes. Every conflict-serializable schedule is view serializable; a view-serializable schedule that is not conflict-serializable has \textbf{blind writes}.
\item
  Number of serial schedules of n transactions: \textbf{n!}.
\end{itemize}

\begin{figure}[!htbp]
\centering
\begin{adjustbox}{max width=\linewidth}
\begin{tikzpicture}
  \node[font=\small] at (0,0) {\begin{tabular}{c|c|c}
     \textbf{T1} & \textbf{T2} & \textbf{T3} \\\hline
     $r_1(A)$ & & \\
     & $w_2(A)$ & \\
     $w_1(A)$ & & \\
     & & $r_3(A)$ \\
  \end{tabular}};
  \begin{scope}[shift={(5.2,0)}]
    \node[vertex] (t1) at (0,0.9) {T1}; \node[vertex] (t2) at (2.2,0.9) {T2}; \node[vertex] (t3) at (1.1,-1.0) {T3};
    \draw[arr, fwine] (t1) to[bend left=22] (t2);
    \draw[arr, fwine] (t2) to[bend left=22] (t1);
    \draw[arr] (t1) -- (t3); \draw[arr] (t2) -- (t3);
  \end{scope}
\end{tikzpicture}
\end{adjustbox}
\caption{A schedule and its precedence graph. The cycle T1 $\leftrightarrow$ T2 (from $r_1(A)$ before $w_2(A)$ and
$w_2(A)$ before $w_1(A)$) means the schedule is not conflict serializable.}
\end{figure}


\needspace{95pt}

\hypertarget{p2-04--recoverability-hierarchy}{%
\subsection{Recoverability hierarchy}\label{p2-04--recoverability-hierarchy}}

\begin{figure}[!htbp]
\centering
\begin{adjustbox}{max width=\linewidth}
\begin{tikzpicture}
  \foreach \w/\h/\col/\name [count=\k] in {13.0/3.6/fgrey/{all schedules}, 10.6/2.85/fslate/recoverable, 8.2/2.1/fsage/{cascadeless (ACA)}, 5.8/1.35/fsand/strict, 3.0/0.6/frose/serial} {
    \filldraw[fill=\col, draw=fgink, line width=0.5pt] (0,0) ellipse ({\w/2} and \h);
    \ifnum\k<5 \node[font=\small] at (0,{\h-0.32}) {\name}; \else \node[font=\small] at (0,0) {\name}; \fi
  }
\end{tikzpicture}
\end{adjustbox}
\caption{Classes of schedules: serial $\subset$ strict $\subset$ cascadeless $\subset$ recoverable.}
\end{figure}


\begin{itemize}
\tightlist
\item
  \textbf{Recoverable}: if Tj reads from Ti, Ti commits before Tj commits.
\item
  \textbf{Cascadeless}: read only committed data.
\item
  \textbf{Strict}: no read/write of item until the last writer commits/aborts.
\end{itemize}

\needspace{196pt}

\hypertarget{p2-04--concurrency-problems}{%
\subsection{Concurrency problems}\label{p2-04--concurrency-problems}}

Lost update (WW), dirty read (WR --- uncommitted dependency), unrepeatable read (RW), phantom read (insertions).
SQL isolation levels: READ UNCOMMITTED \textless{} READ COMMITTED \textless{} REPEATABLE READ \textless{} SERIALIZABLE.

\needspace{130pt}

\hypertarget{p2-04--concurrency-control-protocols}{%
\subsection{Concurrency Control Protocols}\label{p2-04--concurrency-control-protocols}}

\begingroup\small

\par\medskip\noindent\hfil\begin{tabular}{@{}
  >{\raggedright\arraybackslash}p{(\columnwidth - 2\tabcolsep) * \real{0.2161}}
  >{\raggedright\arraybackslash}p{(\columnwidth - 2\tabcolsep) * \real{0.7839}}@{}}
\toprule\noalign{}
\begin{minipage}[b]{\linewidth}\raggedright
\textbf{Protocol}
\end{minipage} & \begin{minipage}[b]{\linewidth}\raggedright
\textbf{Key points}
\end{minipage} \\
\midrule\noalign{}
Lock-based (S/X) & S compatible with S only \\
\textbf{Basic 2PL} & Growing phase then shrinking phase; ensures conflict serializability; \textbf{deadlock possible}, cascading rollback possible \\
\textbf{Strict 2PL} & Hold \textbf{X locks} till commit → cascadeless, strict \\
\textbf{Rigorous 2PL} & Hold \textbf{all locks} till commit; serial order = commit order \\
Conservative 2PL & Acquire all locks before starting → \textbf{deadlock-free} \\
\textbf{Timestamp ordering} & Older first; Read: if TS(T) \textless{} W-TS(X) → rollback; Write: if TS(T) \textless{} R-TS(X) or \textless{} W-TS(X) → rollback. Deadlock-free, may starve \\
\textbf{Thomas\textquotesingle{} write rule} & Ignore obsolete write instead of rollback → allows some view-serializable schedules \\
Validation (optimistic) & Read → validate → write phases \\
Multiversion (MVCC) & Readers don\textquotesingle t block writers \\
Multiple granularity & Intention locks IS, IX, SIX \\
\bottomrule\noalign{}
\end{tabular}\hfil\null\par\medskip


\endgroup

Deadlock prevention with timestamps:

\begin{center}
\renewcommand{\arraystretch}{1.35}
\begin{tabularx}{\linewidth}{@{}l X X@{}}
\toprule
\textbf{Scheme} & \textbf{Older asks for younger's lock} & \textbf{Younger asks for older's lock} \\
\midrule
Wait–die (non-pre-emptive) & older \emph{waits} & younger \emph{dies} (rolled back) \\
Wound–wait (pre-emptive) & older \emph{wounds} (aborts) younger & younger \emph{waits} \\
\bottomrule
\end{tabularx}
\end{center}


\needspace{95pt}

\hypertarget{p2-04--recovery}{%
\subsection{Recovery}\label{p2-04--recovery}}

\begin{itemize}
\tightlist
\item
  \textbf{Log-based}: \texttt{\textless{}T\ start\textgreater{}}, \texttt{\textless{}T,\ X,\ old,\ new\textgreater{}}, \texttt{\textless{}T\ commit\textgreater{}}.
\item
  \textbf{Deferred update} (NO-UNDO/REDO), \textbf{Immediate update} (UNDO/REDO).
\item
  \textbf{Checkpoint} reduces log scanning. After crash: transactions committed after checkpoint → REDO; uncommitted → UNDO.
\item
  \textbf{ARIES} (WAL, LSN, repeating history): Analysis → Redo → Undo.
\item
  \textbf{WAL (Write-Ahead Logging)}: log record written to stable storage before data page.
\item
  Shadow paging: no log needed; copy-on-write page table.
\end{itemize}

\needspace{160pt}

\hypertarget{p2-04--8-file-organisation--indexing}{%
\section{File Organisation \& Indexing}\label{p2-04--8-file-organisation--indexing}}

An index is to a table what the index of a book is to its pages: a small, ordered structure that leads
straight to the wanted rows. B\(^{+}\)-trees dominate because they stay balanced and short even for very large
files.

\begingroup\small

\par\medskip\noindent\hfil\begin{tabular}{@{}
  >{\raggedright\arraybackslash}p{(\columnwidth - 2\tabcolsep) * \real{0.1050}}
  >{\raggedright\arraybackslash}p{(\columnwidth - 2\tabcolsep) * \real{0.4356}}@{}}
\toprule\noalign{}
\begin{minipage}[b]{\linewidth}\raggedright
\textbf{Index}
\end{minipage} & \begin{minipage}[b]{\linewidth}\raggedright
\textbf{Description}
\end{minipage} \\
\midrule\noalign{}
Primary & On ordering \textbf{key} field; sparse (one entry per block) \\
Clustering & On ordering \textbf{non-key} field \\
Secondary & On non-ordering field; \textbf{dense} \\
Dense & Entry for every record \\
Sparse & Entry for some records (needs sorted file) \\
Multilevel & Index on index --- log\_fo(blocks) levels \\
\bottomrule\noalign{}
\end{tabular}\hfil\null\par\medskip


\endgroup

\begin{itemize}
\tightlist
\item
  A file can have \textbf{at most one primary or clustering index} but many secondary indexes.
\end{itemize}

\needspace{95pt}

\hypertarget{p2-04--b-tree-vs-b-tree}{%
\subsection{B-Tree vs B+ Tree}\label{p2-04--b-tree-vs-b-tree}}

\begin{figure}[!htbp]
\centering
\begin{adjustbox}{max width=\linewidth}
\begin{tikzpicture}[kn/.style={rectangle split, rectangle split horizontal, rectangle split parts=#1, draw=fgink, line width=0.5pt,
    fill=fpaper, font=\small, inner sep=3pt, minimum height=6.5mm,
    blur shadow={shadow blur steps=6, shadow xshift=0.7pt, shadow yshift=-0.9pt, shadow opacity=22}}]
  \node[kn=2, fill=fslate] (root) at (0,0) {30 \nodepart{two} 60};
  \node[kn=2, fill=fsage] (l1) at (-4.2,-1.8) {10 \nodepart{two} 20};
  \node[kn=2, fill=fsage] (l2) at (0,-1.8) {30 \nodepart{two} 40};
  \node[kn=3, fill=fsage] (l3) at (4.4,-1.8) {60 \nodepart{two} 70 \nodepart{three} 80};
  \draw[arr] (root.south west) -- (l1.north);
  \draw[arr] (root.south) -- (l2.north);
  \draw[arr] (root.south east) -- (l3.north);
  \draw[arr, fwine] (l1.east) -- (l2.west); \draw[arr, fwine] (l2.east) -- (l3.west);
  \node[lbl, anchor=west] at (1.6,0) {internal node: keys only (guides the search)};
  \node[lbl, text=fwine] at (0,-2.6) {leaves hold every key and are chained for range queries};
\end{tikzpicture}
\end{adjustbox}
\caption{A B$^{+}$-tree of order 3. Unlike a B-tree, internal keys are repeated in the leaves, and the leaf
chain supports sequential access.}
\end{figure}


\begingroup\small

\par\medskip\noindent\hfil\begin{tabular}{@{}
  >{\raggedright\arraybackslash}p{(\columnwidth - 2\tabcolsep) * \real{0.2311}}
  >{\raggedright\arraybackslash}p{(\columnwidth - 2\tabcolsep) * \real{0.4772}}@{}}
\toprule\noalign{}
\begin{minipage}[b]{\linewidth}\raggedright
\textbf{B-Tree}
\end{minipage} & \begin{minipage}[b]{\linewidth}\raggedright
\textbf{B+ Tree}
\end{minipage} \\
\midrule\noalign{}
Data pointers in all nodes & Data pointers only at leaves \\
No duplicate keys & Internal keys repeated in leaves \\
Leaves not linked & Leaves linked → efficient range/\allowbreak{}sequential access \\
Lower fan-out & Higher fan-out (internal nodes smaller) → fewer levels \\
\bottomrule\noalign{}
\end{tabular}\hfil\null\par\medskip


\endgroup

Order p (max children): every node ≤ p children; non-root internal ≥ ⌈p/2⌉ children; root ≥ 2.
\textbf{Order calculation}: block 1024 B, key 9 B, block pointer 6 B, record pointer 7 B.

\begin{itemize}
\tightlist
\item
  B+ internal: p·6 + (p−1)·9 ≤ 1024 → 15p ≤ 1033 → \textbf{p = 68}.
\item
  B+ leaf: q(9 + 7) + 6 ≤ 1024 → \textbf{q = 63}.
\item
  B-tree node: p·6 + (p−1)(9 + 7) ≤ 1024 → 22p ≤ 1040 → \textbf{p = 47}.
\end{itemize}

\textbf{Hashing}: static (overflow chains), dynamic --- extendible hashing (directory doubles, global/local depth), linear hashing.

\needspace{95pt}

\hypertarget{p2-04--query-optimisation}{%
\subsection{Query Optimisation}\label{p2-04--query-optimisation}}

\begin{itemize}
\tightlist
\item
  Heuristics: perform \textbf{selection early}, projection early, most restrictive selection first, avoid Cartesian products, replace × + σ with joins.
\item
  Cost-based: estimate block transfers \& seeks. Join algorithms: nested-loop (nr·bs + br), block nested-loop (br·bs + br), indexed, sort-merge, hash join.
\item
  Query tree; equivalence rules (σ commutes, cascade of π, joins associative/commutative).
\end{itemize}

\needspace{130pt}

\hypertarget{p2-04--9-advanced--enhanced-databases}{%
\section{Advanced \& Enhanced Databases}\label{p2-04--9-advanced--enhanced-databases}}

\begingroup\small

\par\medskip\noindent\hfil\begin{tabular}{@{}
  >{\raggedright\arraybackslash}p{(\columnwidth - 2\tabcolsep) * \real{0.1483}}
  >{\raggedright\arraybackslash}p{(\columnwidth - 2\tabcolsep) * \real{0.8517}}@{}}
\toprule\noalign{}
\begin{minipage}[b]{\linewidth}\raggedright
\textbf{Model}
\end{minipage} & \begin{minipage}[b]{\linewidth}\raggedright
\textbf{Notes}
\end{minipage} \\
\midrule\noalign{}
OODBMS & Objects, OID, inheritance; ODMG standard, OQL \\
ORDBMS & Relational + UDTs, inheritance (PostgreSQL, Oracle) \\
Temporal & Valid time \& transaction time \\
Multimedia & Images, audio, video; content-based retrieval \\
Deductive & Facts + rules; \textbf{Datalog} (Prolog-like) \\
XML DB & XPath, XQuery; native XML DBs \\
Mobile & Intermittent connectivity, caching \\
GIS & Spatial data: raster/\allowbreak{}vector; R-trees, quad-trees \\
Genome & Biological sequence data \\
\textbf{Distributed DB} & Fragmentation (horizontal, vertical, mixed), replication, transparency; \textbf{2PC} (two-phase commit: prepare → commit) --- blocking; 3PC non-blocking \\
\bottomrule\noalign{}
\end{tabular}\hfil\null\par\medskip


\endgroup

\textbf{Security}: discretionary (GRANT/REVOKE, \texttt{WITH\ GRANT\ OPTION}), mandatory (Bell-LaPadula: no read up, no write down), role-based access control, statistical DB inference control, encryption.

\needspace{125pt}

\hypertarget{p2-04--10-data-warehousing--mining}{%
\section{Data Warehousing \& Mining}\label{p2-04--10-data-warehousing--mining}}

Operational databases record today\textquotesingle s transactions; a data warehouse collects years of history from many
sources for analysis. Data mining then searches that history for patterns.

\begin{figure}[!htbp]
\centering
\begin{adjustbox}{max width=\linewidth}
\begin{tikzpicture}
  \node[db, text width=16mm, font=\footnotesize] (s1) at (0,1.0) {OLTP\\sources};
  \node[db, text width=16mm, font=\footnotesize] (s2) at (0,-1.0) {flat\\files};
  \node[sand, text width=20mm, font=\footnotesize] (etl) at (2.9,0) {\textbf{ETL}\\extract,\\transform, load};
  \node[db, minimum width=26mm, minimum height=18mm, text width=24mm, font=\footnotesize] (dw) at (6.1,0) {\textbf{data}\\\textbf{warehouse}};
  \node[sage, text width=18mm, font=\footnotesize] (dm) at (9.1,1.3) {data marts};
  \begin{scope}[shift={(8.4,-1.9)}, x={(1cm,0)}, y={(0,1cm)}, z={(0.45cm,0.3cm)}]
    \foreach \i in {0,1,2} \foreach \j in {0,1,2} {
      \filldraw[fill=fslate!80, draw=fgink, line width=0.35pt] (\i*0.5,\j*0.5,0) -- ++(0.5,0,0) -- ++(0,0.5,0) -- ++(-0.5,0,0) -- cycle;
      \filldraw[fill=fslate!45, draw=fgink, line width=0.35pt] (\i*0.5,1.5,\j*0.5) -- ++(0.5,0,0) -- ++(0,0,0.5) -- ++(-0.5,0,0) -- cycle;
      \filldraw[fill=fslate!95!black, draw=fgink, line width=0.35pt] (1.5,\i*0.5,\j*0.5) -- ++(0,0.5,0) -- ++(0,0,0.5) -- ++(0,-0.5,0) -- cycle;
    }
    \node[font=\scriptsize] at (0.75,-0.25,0) {item};
    \node[font=\scriptsize, rotate=90] at (-0.25,0.75,0) {time};
    \node[font=\scriptsize] at (2.25,0.1,0.9) {location};
  \end{scope}
  \node[font=\footnotesize] at (9.5,-2.85) {OLAP cube};
  \node[rose, text width=22mm, font=\footnotesize] (bi) at (12.4,0) {reports,\\dashboards,\\mining};
  \draw[arr] (s1) -- (etl); \draw[arr] (s2) -- (etl); \draw[arr] (etl) -- (dw);
  \draw[arr] (dw) -- (dm); \draw[arr] (dw) -- (8.2,-1.2);
  \draw[arr] (dm) -- (bi); \draw[arr] (10.6,-1.2) -- (bi);
\end{tikzpicture}
\end{adjustbox}
\caption{Data-warehouse architecture. The warehouse is subject-oriented, integrated, time-variant and
non-volatile (Inmon); OLAP views it as a multidimensional cube that can be sliced, diced, rolled up and drilled down.}
\end{figure}


\begin{itemize}
\tightlist
\item
  Definition by \textbf{W.H. Inmon}: subject-oriented, integrated, time-variant, non-volatile. Kimball: dimensional modelling (bottom-up data marts).
\item
  Schemas: \textbf{Star} (one fact table + denormalised dimensions), \textbf{Snowflake} (normalised dimensions), \textbf{Fact constellation / Galaxy} (multiple fact tables).
\item
  OLAP operations: \textbf{roll-up} (aggregate up / drill-up), \textbf{drill-down}, \textbf{slice} (fix one dimension), \textbf{dice} (sub-cube on multiple), \textbf{pivot} (rotate).
\item
  ROLAP (relational), MOLAP (multidimensional arrays), HOLAP (hybrid).
\item
  Number of cuboids in n-dimensional cube without hierarchies: \textbf{2ⁿ}; with Lᵢ levels: Π(Lᵢ + 1).
\item
  OLTP (current, detailed, many short transactions, normalised) vs OLAP (historical, summarised, complex queries, denormalised).
\end{itemize}

\needspace{155pt}

\hypertarget{p2-04--data-mining-kdd-process}{%
\subsection{Data mining (KDD process)}\label{p2-04--data-mining-kdd-process}}

Selection → Pre-processing → Transformation → \textbf{Data mining} → Interpretation/Evaluation.

\textbf{Association rules} (Apriori):

\begin{formulatable}
Support of $A \Rightarrow B$ & \frac{\text{count}(A \cup B)}{N} \\
Confidence of $A \Rightarrow B$ & \frac{\text{support}(A \cup B)}{\text{support}(A)} \\
Lift & \frac{\text{confidence}(A\Rightarrow B)}{\text{support}(B)}\quad(>1\text{: positive correlation}) \\
Apriori property & \text{every subset of a frequent itemset is frequent} \\
FP-growth & \text{mines an FP-tree without generating candidates} \\
\end{formulatable}


\begin{itemize}
\tightlist
\item
  Classification (supervised): decision tree (ID3 -- information gain, C4.5 -- gain ratio, CART -- Gini), Naive Bayes, k-NN, SVM.
\item
  Clustering (unsupervised): k-means (partitioning), k-medoids (PAM), hierarchical (agglomerative/divisive), DBSCAN (density).
\item
  Outlier detection, regression.
\end{itemize}

\needspace{125pt}

\hypertarget{p2-04--11-big-data--nosql}{%
\section{Big Data \& NoSQL}\label{p2-04--11-big-data--nosql}}

When data become too large, too fast or too varied for a single relational server, systems distribute them
across many machines and often relax strict consistency in exchange for availability.

\begin{itemize}
\tightlist
\item
  \textbf{5 Vs}: Volume, Velocity, Variety, Veracity, Value (3 Vs originally --- Doug Laney).
\item
  Types: structured, semi-structured (JSON, XML), unstructured (text, video).
\item
  \textbf{Hadoop}: HDFS (NameNode = master/metadata, DataNode = storage; default block 128 MB, replication 3) + \textbf{MapReduce} (Map → Shuffle \& Sort → Reduce) + \textbf{YARN} (ResourceManager, NodeManager). Ecosystem: Hive (SQL-like), Pig (Pig Latin), HBase (column store), Sqoop (RDBMS ↔ HDFS), Flume (logs), ZooKeeper (coordination), Oozie (workflow), Spark (in-memory, RDDs).
\item
  \textbf{CAP theorem} (Brewer): a distributed system can guarantee only \textbf{2 of 3}: Consistency, Availability, Partition tolerance. (CP: HBase, MongoDB; AP: Cassandra, CouchDB, DynamoDB.)
\item
  \textbf{BASE}: Basically Available, Soft state, Eventual consistency (vs ACID).
\end{itemize}

\begingroup\small

\par\medskip\noindent\hfil\begin{tabular}{@{}
  >{\raggedright\arraybackslash}p{(\columnwidth - 2\tabcolsep) * \real{0.1450}}
  >{\raggedright\arraybackslash}p{(\columnwidth - 2\tabcolsep) * \real{0.2600}}@{}}
\toprule\noalign{}
\begin{minipage}[b]{\linewidth}\raggedright
\textbf{NoSQL type}
\end{minipage} & \begin{minipage}[b]{\linewidth}\raggedright
\textbf{Examples}
\end{minipage} \\
\midrule\noalign{}
Key--value & Redis, DynamoDB, Riak \\
Document & MongoDB, CouchDB \\
Column-family & Cassandra, HBase, Bigtable \\
Graph & Neo4j, JanusGraph \\
\bottomrule\noalign{}
\end{tabular}\hfil\null\par\medskip


\endgroup

\needspace{196pt}

\hypertarget{p2-04--12-deeper-dive--sql-on-real-tables-minimal-cover-full-normalisation--schedules}{%
\section{Deeper Dive --- SQL on Real Tables, Minimal Cover, Full Normalisation \& Schedules}\label{p2-04--12-deeper-dive--sql-on-real-tables-minimal-cover-full-normalisation--schedules}}

\needspace{160pt}

\hypertarget{p2-04--121-sql-worked-on-sample-data}{%
\subsection{SQL Worked on Sample Data}\label{p2-04--121-sql-worked-on-sample-data}}

\textbf{EMP}

\begingroup\small

\par\medskip\noindent\hfil\begin{tabular}{@{}
  >{\raggedright\arraybackslash}p{(\columnwidth - 6\tabcolsep) * \real{0.0479}}
  >{\raggedright\arraybackslash}p{(\columnwidth - 6\tabcolsep) * \real{0.0720}}
  >{\raggedright\arraybackslash}p{(\columnwidth - 6\tabcolsep) * \real{0.0821}}
  >{\raggedright\arraybackslash}p{(\columnwidth - 6\tabcolsep) * \real{0.0779}}@{}}
\toprule\noalign{}
\begin{minipage}[b]{\linewidth}\raggedright
\textbf{eid}
\end{minipage} & \begin{minipage}[b]{\linewidth}\raggedright
\textbf{name}
\end{minipage} & \begin{minipage}[b]{\linewidth}\raggedright
\textbf{dept}
\end{minipage} & \begin{minipage}[b]{\linewidth}\raggedright
\textbf{salary}
\end{minipage} \\
\midrule\noalign{}
1 & Asha & CS & 50000 \\
2 & Ravi & CS & 40000 \\
3 & Meena & EE & 45000 \\
4 & John & EE & 30000 \\
5 & Priya & ME & 60000 \\
6 & Kiran & NULL & 35000 \\
\bottomrule\noalign{}
\end{tabular}\hfil\null\par\medskip


\endgroup

\textbf{DEPT}

\begingroup\small

\par\medskip\noindent\hfil\begin{tabular}{@{}
  >{\raggedright\arraybackslash}p{(\columnwidth - 2\tabcolsep) * \real{0.0546}}
  >{\raggedright\arraybackslash}p{(\columnwidth - 2\tabcolsep) * \real{0.0907}}@{}}
\toprule\noalign{}
\begin{minipage}[b]{\linewidth}\raggedright
\textbf{dept}
\end{minipage} & \begin{minipage}[b]{\linewidth}\raggedright
\textbf{building}
\end{minipage} \\
\midrule\noalign{}
CS & B1 \\
EE & B2 \\
CE & B3 \\
\bottomrule\noalign{}
\end{tabular}\hfil\null\par\medskip


\endgroup

\begingroup\small

\par\medskip\noindent\hfil\begin{tabular}{@{}
  >{\raggedright\arraybackslash}p{(\columnwidth - 4\tabcolsep) * \real{0.0402}}
  >{\raggedright\arraybackslash}p{(\columnwidth - 4\tabcolsep) * \real{0.6071}}
  >{\raggedright\arraybackslash}p{(\columnwidth - 4\tabcolsep) * \real{0.3527}}@{}}
\toprule\noalign{}
\begin{minipage}[b]{\linewidth}\raggedright
\textbf{\#}
\end{minipage} & \begin{minipage}[b]{\linewidth}\raggedright
\textbf{Query}
\end{minipage} & \begin{minipage}[b]{\linewidth}\raggedright
\textbf{Result}
\end{minipage} \\
\midrule\noalign{}
1 & \texttt{SELECT\ dept,\ COUNT(*),\ AVG(salary)\ FROM\ EMP\ GROUP\ BY\ dept\ HAVING\ COUNT(*)\ \textgreater{}\ 1} & (CS, 2, 45000), (EE, 2, 37500) \\
2 & \texttt{SELECT\ name\ FROM\ EMP\ WHERE\ salary\ \textgreater{}\ (SELECT\ AVG(salary)\ FROM\ EMP)} & avg = 43333.33 → Asha, Meena, Priya \\
3 & \texttt{SELECT\ COUNT(dept),\ COUNT(DISTINCT\ dept)\ FROM\ EMP} & 5, 3 (NULL ignored) \\
4 & \texttt{EMP\ JOIN\ DEPT\ ON\ EMP.dept\ =\ DEPT.dept} & 4 rows (Asha, Ravi, Meena, John) \\
5 & \texttt{EMP\ LEFT\ JOIN\ DEPT\ …} & 6 rows (Priya \& Kiran get NULL building) \\
6 & \texttt{EMP\ RIGHT\ JOIN\ DEPT\ …} & 5 rows (CE appears with NULL employee) \\
7 & \texttt{EMP\ FULL\ OUTER\ JOIN\ DEPT\ …} & 7 rows \\
8 & Correlated: \texttt{SELECT\ name\ FROM\ EMP\ e1\ WHERE\ salary\ =\ (SELECT\ MAX(salary)\ FROM\ EMP\ e2\ WHERE\ e2.dept\ =\ e1.dept)} & Asha, Meena, Priya (Kiran excluded: NULL = NULL is UNKNOWN) \\
9 & \texttt{SELECT\ dept\ FROM\ DEPT\ WHERE\ dept\ NOT\ IN\ (SELECT\ dept\ FROM\ EMP)} & \textbf{Empty!} The subquery contains NULL, so every NOT IN test is UNKNOWN \\
10 & Same with \texttt{NOT\ EXISTS\ (SELECT\ *\ FROM\ EMP\ WHERE\ EMP.dept\ =\ DEPT.dept)} & CE ✔ --- NOT EXISTS is NULL-safe \\
\bottomrule\noalign{}
\end{tabular}\hfil\null\par\medskip


\endgroup

\needspace{125pt}

\hypertarget{p2-04--122-minimal-canonical-cover--worked}{%
\subsection{Minimal (Canonical) Cover --- Worked}\label{p2-04--122-minimal-canonical-cover--worked}}

F = \{A → BC, B → C, A → B, AB → C\}

\begin{center}
\renewcommand{\arraystretch}{1.35}
\begin{tabularx}{\linewidth}{@{}l >{\raggedright\arraybackslash}X@{}}
\toprule
\textbf{Step} & \textbf{Result} \\
\midrule
1. Split right-hand sides & $A\to B,\ A\to C,\ B\to C,\ A\to B,\ AB\to C$ \\
2. Remove duplicates & $A\to B,\ A\to C,\ B\to C,\ AB\to C$ \\
3. Remove extraneous LHS attributes & in $AB\to C$, $B$ is extraneous since $A^{+}\ni C$; gives $A\to C$ (duplicate) \\
4. Remove redundant FDs & $A\to C$ follows from $A\to B$ and $B\to C$ \\
\midrule
Minimal cover & $F_c = \{A\to B,\ B\to C\}$ \\
\bottomrule
\end{tabularx}
\end{center}


\needspace{125pt}

\hypertarget{p2-04--123-normalisation--from-1nf-to-bcnf}{%
\subsection{Normalisation --- From 1NF to BCNF}\label{p2-04--123-normalisation--from-1nf-to-bcnf}}

R(S, C, I, P, G) --- Student, Course, Instructor, instructor Phone, Grade.
F = \{ SC → G, C → I, I → P \}. S and C never appear on the right → \textbf{candidate key = SC}.

\begin{figure}[!htbp]
\centering
\begin{adjustbox}{max width=\linewidth}
\begin{tikzpicture}
  \node[grey, text width=38mm] (R) at (0,0) {$R(S, C, I, P, G)$\\\footnotesize key $SC$ — 1NF only};
  \node[sage, text width=26mm] (R1) at (-3.4,-1.8) {$R_1(S, C, G)$\\\footnotesize BCNF};
  \node[slate, text width=32mm] (R2) at (3.2,-1.8) {$R_2(C, I, P)$\\\footnotesize 2NF; $C\to I\to P$};
  \node[sage, text width=22mm] (R21) at (1.6,-3.6) {$R_{21}(C, I)$\\\footnotesize BCNF};
  \node[sage, text width=22mm] (R22) at (4.8,-3.6) {$R_{22}(I, P)$\\\footnotesize BCNF};
  \draw[arr] (R) -- node[lbl, above left] {remove partial $C\to I, P$} (R1);
  \draw[arr] (R) -- (R2);
  \draw[arr] (R2) -- node[lbl, left] {remove transitive} (R21); \draw[arr] (R2) -- (R22);
\end{tikzpicture}
\end{adjustbox}
\caption{Decomposing $R(S,C,I,P,G)$ with $SC\to G$, $C\to I$, $I\to P$. The final schemas are lossless and
preserve every dependency.}
\end{figure}


Final schema: \textbf{(S, C, G), (C, I), (I, P)} --- every determinant is a key → BCNF; decomposition is lossless (common attributes C and I are keys of a side) and preserves all three FDs.

\needspace{130pt}

\hypertarget{p2-04--124-recoverability-examples-w--write-r--read-c--commit}{%
\subsection{Recoverability Examples (W = write, R = read, C = commit)}\label{p2-04--124-recoverability-examples-w--write-r--read-c--commit}}

\begingroup\small

\par\medskip\noindent\hfil\begin{tabular}{@{}
  >{\raggedright\arraybackslash}p{(\columnwidth - 2\tabcolsep) * \real{0.2050}}
  >{\raggedright\arraybackslash}p{(\columnwidth - 2\tabcolsep) * \real{0.7661}}@{}}
\toprule\noalign{}
\begin{minipage}[b]{\linewidth}\raggedright
\textbf{Schedule}
\end{minipage} & \begin{minipage}[b]{\linewidth}\raggedright
\textbf{Classification}
\end{minipage} \\
\midrule\noalign{}
W1(A) R2(A) C2 C1 & \textbf{Not recoverable} --- T2 commits after reading T1\textquotesingle s uncommitted data, before T1 commits \\
W1(A) R2(A) C1 C2 & Recoverable, but not cascadeless (abort of T1 forces abort of T2) \\
W1(A) C1 R2(A) C2 & Cascadeless (reads only committed data) \\
W1(A) W2(A) C1 C2 & Cascadeless but \textbf{not strict} (T2 overwrote uncommitted A) \\
\bottomrule\noalign{}
\end{tabular}\hfil\null\par\medskip


\endgroup

\needspace{125pt}

\hypertarget{p2-04--125-precedence-graphs}{%
\subsection{Precedence Graphs}\label{p2-04--125-precedence-graphs}}

Schedule S: R1(A) W1(A) R2(A) W2(A) R1(B) W1(B) R2(B) W2(B)

\begin{figure}[!htbp]
\centering
\begin{adjustbox}{max width=\linewidth}
\begin{tikzpicture}[vertex/.append style={minimum size=10mm, font=\normalsize}]
  \node[vertex] (t1) at (0,0) {T1}; \node[vertex] (t2) at (5,0) {T2};
  \draw[arr] (t1) -- node[lbl, above, align=center] {$w_1(A)$ before $r_2(A)$\\$w_1(B)$ before $r_2(B)$} (t2);
  \node[capt] at (2.5,-1.6) {(a) acyclic: equivalent to T1 $\to$ T2};
  \begin{scope}[shift={(8,0)}]
    \node[vertex] (u1) at (0,0) {T1}; \node[vertex] (u2) at (5,0) {T2};
    \draw[arr, fwine] (u1) to[bend left=30] node[lbl, above] {$r_1(X)$ before $w_2(X)$} (u2);
    \draw[arr, fwine] (u2) to[bend left=30] node[lbl, below] {$r_2(Y)$ before $w_1(Y)$} (u1);
    \node[capt] at (2.5,-1.6) {(b) cycle: not conflict serializable};
  \end{scope}
\end{tikzpicture}
\end{adjustbox}
\caption{Precedence graphs for the two schedules S and S$'$ below.}
\end{figure}


Acyclic → conflict serializable, equivalent to serial order \textbf{T1 → T2}.

Schedule S′: R1(X) R2(Y) W2(X) W1(Y)

% The cyclic graph is drawn beside the acyclic one in p2-04-prec-ok.


Cycle → \textbf{not} conflict serializable.

\needspace{95pt}

\hypertarget{p2-04--126-two-phase-locking--lock-point}{%
\subsection{Two-Phase Locking --- Lock Point}\label{p2-04--126-two-phase-locking--lock-point}}

\begin{figure}[!htbp]
\centering
\begin{adjustbox}{max width=\linewidth}
\begin{tikzpicture}
  \draw[arr] (0,0) -- (10.5,0) node[right, font=\footnotesize] {time};
  \draw[arr] (0,0) -- (0,3.4) node[above, font=\footnotesize] {locks held};
  \fill[fsage!70] (0.5,0) -- (1.5,1) -- (2.5,1) -- (2.5,2) -- (3.6,2) -- (3.6,3) -- (5.2,3) -- (5.2,2) -- (6.6,2) -- (6.6,1) -- (8.2,1) -- (8.2,0) -- cycle;
  \draw[fgink, line width=0.6pt] (0.5,0) -- (1.5,1) -- (2.5,1) -- (2.5,2) -- (3.6,2) -- (3.6,3) -- (5.2,3) -- (5.2,2) -- (6.6,2) -- (6.6,1) -- (8.2,1) -- (8.2,0);
  \fill[fwine] (3.6,3) circle (2.4pt);
  \node[font=\footnotesize, text=fwine, above] at (3.6,3.05) {lock point};
  \node[lbl] at (2.0,2.4) {growing phase};
  \node[lbl] at (7.2,2.4) {shrinking phase};
\end{tikzpicture}
\end{adjustbox}
\caption{Two-phase locking: a transaction acquires all its locks before releasing any. Serializability order
follows the order of the lock points.}
\end{figure}


\needspace{95pt}

\hypertarget{p2-04--127-extendible-hashing-sketch}{%
\subsection{Extendible Hashing (sketch)}\label{p2-04--127-extendible-hashing-sketch}}

\begin{figure}[!htbp]
\centering
\begin{adjustbox}{max width=\linewidth}
\begin{tikzpicture}[cell/.style={draw=fgink, line width=0.5pt, fill=fslate, minimum width=11mm, minimum height=7mm, font=\small\ttfamily},
    bkt/.style={rectangle split, rectangle split horizontal, rectangle split parts=3, draw=fgink, line width=0.5pt, fill=fsand, font=\small, minimum height=7mm, inner sep=3pt}]
  \node[cell] (d00) at (0,1.05) {00}; \node[cell] (d01) at (0,0.35) {01};
  \node[cell] (d10) at (0,-0.35) {10}; \node[cell] (d11) at (0,-1.05) {11};
  \node[capt] at (0,1.85) {directory (global depth 2)};
  \node[bkt] (b1) at (5,1.4) {4 \nodepart{two} 12 \nodepart{three} 32};
  \node[bkt] (b2) at (5,0) {1 \nodepart{two} 5 \nodepart{three} 21};
  \node[bkt, rectangle split parts=2] (b3) at (5,-1.4) {10 \nodepart{two} 6};
  \node[lbl, anchor=west] at (6.6,1.4) {local depth 2}; \node[lbl, anchor=west] at (6.6,0) {local depth 1 — shared by 01 and 11}; \node[lbl, anchor=west] at (6.2,-1.4) {local depth 2};
  \draw[arr] (d00.east) -- (b1.west); \draw[arr] (d01.east) -- (b2.west); \draw[arr] (d11.east) -- (b2.west); \draw[arr] (d10.east) -- (b3.west);
\end{tikzpicture}
\end{adjustbox}
\caption{Extendible hashing. An overflowing bucket whose local depth equals the global depth forces the
directory to double; otherwise only the bucket splits.}
\end{figure}


\needspace{95pt}

\hypertarget{p2-04--128-data-cube-lattice-3-dimensions--2uxb3--8-cuboids}{%
\subsection{Data Cube Lattice (3 dimensions → 2³ = 8 cuboids)}\label{p2-04--128-data-cube-lattice-3-dimensions--2uxb3--8-cuboids}}

\begin{figure}[!htbp]
\centering
\begin{adjustbox}{max width=\linewidth}
\begin{tikzpicture}[c/.style={box, fill=fslate, font=\footnotesize, minimum width=26mm}]
  \node[c, fill=fsand] (A) at (0,3) {(time, item, location)};
  \node[c] (B) at (-4,1.5) {(time, item)}; \node[c] (C) at (0,1.5) {(time, location)}; \node[c] (D) at (4,1.5) {(item, location)};
  \node[c, fill=fsage] (E) at (-4,0) {(time)}; \node[c, fill=fsage] (F) at (0,0) {(item)}; \node[c, fill=fsage] (G) at (4,0) {(location)};
  \node[c, fill=frose] (H) at (0,-1.5) {( ) — apex};
  \foreach \a/\b in {A/B, A/C, A/D, B/E, B/F, C/E, C/G, D/F, D/G, E/H, F/H, G/H} \draw[thinline] (\a) -- (\b);
  \node[lbl, anchor=west] at (5.6,3) {base cuboid (most detail)};
  \node[lbl, anchor=west] at (5.6,-1.5) {grand total};
\end{tikzpicture}
\end{adjustbox}
\caption{The lattice of $2^3 = 8$ cuboids for three dimensions. Rolling up moves down the lattice;
drilling down moves up.}
\end{figure}


\needspace{125pt}

\hypertarget{p2-04--previous-year-questions-pyq-pattern}{%
\section{Previous Year Questions (PYQ pattern)}\label{p2-04--previous-year-questions-pyq-pattern}}

\textbf{Architecture \& ER}

\begin{enumerate}
\def\labelenumi{\arabic{enumi}.}
\tightlist
\item
  Ability to change the conceptual schema without changing application programs: \textbf{logical data independence}
\item
  The overall design of the database is called: \textbf{schema}
\item
  A weak entity set is represented by: \textbf{double rectangle}
\item
  Minimum number of tables for E1 (1)--R--(N) E2 with no attributes on R: \textbf{2}
\item
  Minimum tables for M:N relationship between two entities: \textbf{3}
\item
  Which is a DCL command? \textbf{GRANT}
\item
  TRUNCATE is a: \textbf{DDL command}
\end{enumerate}

\textbf{Keys}

\begin{enumerate}
\def\labelenumi{\arabic{enumi}.}
\setcounter{enumi}{7}
\tightlist
\item
  R(A,B,C,D) with only candidate key A --- number of super keys: 2³ = \textbf{8}
\item
  R(A,B,C,D) with candidate keys A and B: 8 + 8 − 4 = \textbf{12}
\item
  Which key cannot be NULL? \textbf{Primary key}
\end{enumerate}

\textbf{Relational algebra / SQL}

\begin{enumerate}
\def\labelenumi{\arabic{enumi}.}
\setcounter{enumi}{10}
\tightlist
\item
  R has 10 tuples, S has 5; R × S has: \textbf{50}
\item
  Which RA operation answers "for all" queries? \textbf{Division}
\item
  Which is not a fundamental RA operation? (a) σ (b) π (c) ∩ (d) − --- \textbf{Ans: (c)}
\item
  \texttt{SELECT\ COUNT(*)} vs \texttt{COUNT(salary)} on a table with 10 rows and 2 NULL salaries: \textbf{10 and 8}
\item
  HAVING is applied: \textbf{after GROUP BY on groups}
\item
  \texttt{salary\ \textgreater{}\ ALL\ (SELECT\ salary\ FROM\ emp\ WHERE\ dept\ =\ 5)} means: \textbf{greater than the maximum salary of dept 5}
\item
  Which set operator retains duplicates? \textbf{UNION ALL}
\item
  A trigger follows which model? \textbf{Event--Condition--Action}
\item
  Best protection against SQL injection: \textbf{parameterised/prepared statements}
\end{enumerate}

\textbf{Normalisation}

\begin{enumerate}
\def\labelenumi{\arabic{enumi}.}
\setcounter{enumi}{19}
\tightlist
\item
  R(A,B,C,D), F = \{A→B, B→C, C→D\}. Candidate key? \textbf{A}; Normal form? \textbf{2NF} (transitive dependencies exist)
\item
  R(A,B,C,D), F = \{AB→C, C→D\}. Key AB; C→D is transitive → \textbf{2NF}, not 3NF
\item
  R(A,B,C), F = \{AB→C, C→A\}. Keys AB, CB. Highest NF: \textbf{3NF} (C→A, A is prime)
\item
  Every relation with two attributes is in: \textbf{BCNF}
\item
  Decomposition of R(A,B,C) with A→B into (A,B),(A,C) is: \textbf{lossless} (common A → AB)
\item
  Which NF deals with multivalued dependency? \textbf{4NF}; join dependency? \textbf{5NF}
\item
  BCNF decomposition always guarantees: \textbf{lossless join} (not dependency preservation)
\item
  Closure of \{A\} under \{A→B, B→C, C→D, D→E\}: \textbf{ABCDE}
\end{enumerate}

\textbf{Transactions}

\begin{enumerate}
\def\labelenumi{\arabic{enumi}.}
\setcounter{enumi}{27}
\tightlist
\item
  Which ACID property is ensured by the concurrency control manager? \textbf{Isolation}
\item
  A schedule is conflict serializable iff its precedence graph is: \textbf{acyclic}
\item
  2PL ensures: \textbf{conflict serializability} (not freedom from deadlock)
\item
  Strict 2PL avoids: \textbf{cascading rollbacks}
\item
  In wait-die, if an older transaction requests a lock held by a younger: \textbf{older waits}
\item
  In wound-wait, if an older transaction requests: \textbf{younger is aborted (wounded)}
\item
  Thomas\textquotesingle{} write rule modifies: \textbf{timestamp ordering protocol}
\item
  Number of serial schedules for 4 transactions: \textbf{24}
\item
  Reading uncommitted data is: \textbf{dirty read}
\item
  Write-ahead logging ensures: \textbf{log record on stable storage before data item}
\item
  After a crash, a transaction that has \texttt{\textless{}start\textgreater{}} and \texttt{\textless{}commit\textgreater{}} in log is: \textbf{redone}
\end{enumerate}

\textbf{Indexing}

\begin{enumerate}
\def\labelenumi{\arabic{enumi}.}
\setcounter{enumi}{38}
\tightlist
\item
  B+ tree leaves are linked to support: \textbf{range queries / sequential access}
\item
  Max number of primary indexes on a file: \textbf{1}
\item
  Order of B+ tree internal node: block 512 B, key 10 B, pointer 6 B: p·6 + (p−1)·10 ≤ 512 → 16p ≤ 522 → \textbf{p = 32}
\item
  Secondary index is usually: \textbf{dense}
\end{enumerate}

\textbf{Warehousing / mining / big data}

\begin{enumerate}
\def\labelenumi{\arabic{enumi}.}
\setcounter{enumi}{42}
\tightlist
\item
  Data warehouse is NOT: (a) subject-oriented (b) volatile (c) time-variant (d) integrated --- \textbf{Ans: (b)}
\item
  OLAP operation that rotates the axes: \textbf{pivot}
\item
  Star schema has: \textbf{one fact table, denormalised dimension tables}
\item
  Support of \{bread, butter\} in 5 of 20 transactions: \textbf{25\%}; if bread in 10 → confidence(bread→butter) = \textbf{50\%}
\item
  Apriori uses the property: \textbf{all subsets of a frequent itemset are frequent}
\item
  CAP theorem: choose \textbf{2 of Consistency, Availability, Partition tolerance}
\item
  MongoDB is a: \textbf{document store}; Cassandra: \textbf{column-family}; Neo4j: \textbf{graph}
\item
  In Hadoop, metadata is stored by: \textbf{NameNode}
\item
  K-means is a: \textbf{partitioning clustering} algorithm
\end{enumerate}

\textbf{More practice questions}

\begin{enumerate}
\def\labelenumi{\arabic{enumi}.}
\setcounter{enumi}{51}
\tightlist
\item
  A NOT IN subquery whose result contains NULL returns: \textbf{no rows}
\item
  COUNT(DISTINCT dept) on values \{CS, CS, EE, NULL\}: \textbf{2}
\item
  Minimal cover of \{A → B, B → C, A → C\}: \textbf{\{A → B, B → C\}}
\item
  R(A, B, C, D), F = \{A → B, C → D\}. Candidate key: \textbf{AC}; highest normal form: \textbf{1NF} (partial dependencies)
\item
  A schedule in which transactions read only committed data is: \textbf{cascadeless}
\item
  Schedule R1(X) R2(Y) W2(X) W1(Y) is: \textbf{not conflict serializable}
\item
  In 2PL, the equivalent serial order follows the: \textbf{lock points}
\item
  In extendible hashing, the directory doubles when an overflowing bucket\textquotesingle s local depth: \textbf{equals the global depth}
\item
  Number of cuboids for 4 dimensions without hierarchies: \textbf{16}
\item
  Apex cuboid represents: \textbf{the grand total (all dimensions aggregated)}
\item
  Full outer join of R (5 tuples) and S (3 tuples) with 2 matching pairs (one-to-one): 2 + 3 + 1 = \textbf{6 tuples}
\end{enumerate}

\begin{revbox}

\begin{itemize}
\tightlist
\item
  3NF: X superkey OR A prime · BCNF: X superkey
\item
  Lossless: (R1∩R2) → R1 or R2 · 3NF = lossless + dep. preserving
\item
  Conflict serializable ⇔ acyclic precedence graph · 2PL → serializable but deadlock possible
\item
  Wait-die: old waits, young dies · Wound-wait: old wounds, young waits
\item
  CAP: pick 2 · HDFS NameNode/DataNode · MapReduce Map→Shuffle→Reduce
\end{itemize}

\end{revbox}
